% Section - References
% Roberto Masocco <roberto.masocco@uniroma2.it>
% April 12, 2026

% ### References ###
\section{References}
\graphicspath{{figs/references/}}

% --- What is a reference? ---
\begin{frame}{References}{What is a reference?}
	\justifying
	A \textbg{reference} is an \textbf{alias} — an alternative name for an already-existing object.\\
	\bigskip
	Declared with \texttt{\&} after the type:
	\begin{itemize}
		\item \texttt{int x = 42;}
		\item \texttt{int\& ref = x;} \quad — \texttt{ref} and \texttt{x} now name the \textbg{same object}.
	\end{itemize}
	\bigskip
	References differ from pointers in three important ways:
	\begin{itemize}
		\item a reference \textbg{cannot be null} — it must be bound to an object at declaration;
		\item a reference \textbg{cannot be rebound} — once initialized, it always refers to the same object;
		\item no explicit dereferencing is needed — syntax is \textbg{identical to using the object directly}.
	\end{itemize}
	\bigskip
	\begin{block}{}
		\centering
		A reference is \textbf{not} a new object. No copy is made, no new memory is allocated.
	\end{block}
\end{frame}

% --- References as function parameters ---
\begin{frame}{References}{References as function parameters}
	\justifying
	The most common use of references is in \textbg{function parameter lists}.\\
	\bigskip
	Passing by value in C copies the argument — expensive for large objects, impossible for non-copyable ones.\\
	\bigskip
	Passing by reference gives the function \textbg{direct access} to the caller's object:
	\begin{itemize}
		\item \textbg{modifiable reference} (\texttt{T\&}): the function may read and \textbf{write} the argument;
		\item \textbg{constant reference} (\texttt{const T\&}): the function may only \textbf{read} the argument — no copy, no modification.
	\end{itemize}
	\bigskip
	\texttt{const T\&} is the idiomatic way to pass \textbg{large objects} that should not be modified:
	\begin{itemize}
		\item avoids the cost of a copy;
		\item the \texttt{const} qualifier makes the intent explicit and is enforced by the compiler.
	\end{itemize}
\end{frame}

% --- References vs pointers ---
\begin{frame}{References}{References vs pointers}
	\justifying
	Both refer to an object without copying it, but the choice between them conveys \textbg{intent}.\\
	\bigskip
	Prefer a \textbg{reference} when:
	\begin{itemize}
		\item the referred object \textbg{always exists} for the lifetime of the reference;
		\item no rebinding or pointer arithmetic is needed;
		\item the code should read as naturally as if using the object directly.
	\end{itemize}
	\bigskip
	Prefer a \textbg{pointer} when:
	\begin{itemize}
		\item the object may be \textbg{absent} (\texttt{nullptr} is a valid state);
		\item ownership or lifetime \textbg{must be managed} explicitly;
		\item pointer arithmetic or array indexing is required.
	\end{itemize}
	\bigskip
	\begin{block}{}
		\centering
		In modern C++, raw pointers are used \textbf{only for non-owning observation}.\\
		Ownership is expressed through \textbf{smart pointers} — more on this in the next lecture.
	\end{block}
\end{frame}
